\documentclass[12pt]{article}
\usepackage[T1]{fontenc}
\usepackage[utf8]{inputenc}
\usepackage{lmodern}
\usepackage{textcomp}
\usepackage{enumitem}
\usepackage{geometry}
\usepackage{graphicx}
\usepackage{amsmath, amssymb}
\geometry{margin=1in}
\begin{document}
\begin{center}
Economics - Undergraduate Level Assessment 12 \
Total Marks: 40 \
Answer all questions.
\end{center}

\section*{Multiple Choice (10 marks)}
\begin{enumerate}[label=\arabic*.]
\item Suppose the government seeks to achieve a balanced budgetby levying taxes of $100 billion and making expendituresof $100 billion. How will this affect NNP if MPC = 0.8?
\begin{enumerate}[label=(\alph*)]
    \item Decrease of \$500 billion in NNP
    \item Increase of \$500 billion in NNP
    \item Decrease of \$100 billion in NNP
    \item No change in NNP because taxes cancel out expenditures
    \item Increase of \$200 billion in NNP
\end{enumerate}
\item In what ways is contractionary fiscal policy in the United States likely to affect domestic interest rates and the international value of the dollar?
\begin{enumerate}[label=(\alph*)]
    \item Interest rates increase and the dollar depreciates.
    \item Interest rates remain the same and the dollar depreciates.
    \item Interest rates decrease and the dollar appreciates.
    \item Interest rates and the value of the dollar both remain the same.
    \item Interest rates increase but the value of the dollar remains the same.
\end{enumerate}
\item For a competitive firm, what is the most important thing to consider in deciding whether to shut down in the short run?
\begin{enumerate}[label=(\alph*)]
    \item Do not produce if the TFC is not covered by revenue.
    \item Do not produce if the price is lower than AVC.
    \item Compare AVC to MR.
    \item Compare AFC to MR.
    \item Shut down if the firm is not making a profit.
\end{enumerate}
\item A competitive market for coffee, a normal good, is currently in equilibrium. Which of the following would most likely result in an increase in the demand for coffee?
\begin{enumerate}[label=(\alph*)]
    \item The price of coffee machines rises.
    \item The wage of coffee plantation workers falls.
    \item The price of tea rises.
    \item A decrease in the popularity of energy drinks.
    \item Technology in the harvesting of coffee beans improves.
\end{enumerate}
\item Why are slight differences in growth rate percentages so important?
\begin{enumerate}[label=(\alph*)]
    \item They signify shifts in currency exchange rates
    \item They are indicative of immediate changes in consumer confidence
    \item They determine the exact size of the workforce
    \item They indicate economic stability
    \item They show population change
\end{enumerate}
\end{enumerate}

\section*{True / False (10 marks)}
\textit{Answer True or False}
\begin{enumerate}[label=\arabic*.]
\setcounter{enumi}{5}
\item True or False: The elasticity of supply is typically greater when producers have less capital to invest.
\item True or False: When the U.S. dollar appreciates, U.S. residents are more likely to take vacations in foreign countries due to lower foreign prices.
\item True or False: Wage earners gain from anticipated inflation due to rising wages that outpace inflation, increasing their purchasing power.
\item True or False: Models in economics always lead to increased data accuracy and improved decision-making outcomes for policymakers.
\item True or False: International trade and subsidies competition are the primary reasons for advocating special policies for farmers.
\end{enumerate}

\section*{Long Form Response (20 marks)}
\textit{Answer each question with a clear and complete explanation.}
\begin{enumerate}[label=\arabic*.]
\setcounter{enumi}{10}
\item Discuss the relationship between average product (AP) and marginal product (MP) of labor, explaining how changes in MP influence AP and the implications for production efficiency.
\item Discuss the inefficiencies associated with a monopolistically competitive firm's profit- maximizing equilibrium, particularly focusing on how marginal cost being below price affects societal welfare and output levels.
\end{enumerate}

\vfill
\noindent\textit{End of Paper}
\end{document}